\documentclass[11pt]{article}
\usepackage[margin=1in]{geometry}
\usepackage{booktabs,array,amsmath,xcolor,etoolbox,listings,lineno}
\usepackage[colorlinks=true,linkcolor=blue,urlcolor=blue,citecolor=blue]{hyperref}
\setlength{\tabcolsep}{3pt}
\AtBeginEnvironment{table}{\small}
\sloppy

% Red TODO marker so mentors can see every open item.
\newcommand{\todo}[1]{\textcolor{red}{\textbf{[TODO: #1]}}}
% Tables are generated by make_paper_tables.py into tables/. If a file is missing, show a visible placeholder.
\newcommand{\tablein}[1]{\IfFileExists{tables/#1}{\input{tables/#1}}{\par\todo{table file \detokenize{#1} missing}\par}}

\lstset{basicstyle=\ttfamily\footnotesize,breaklines=true,frame=single,columns=fullflexible}

\title{Repeated-Trial Evaluation of Two LLMs for Express.js Unit Test Generation:\\ Results on Author-Written and External Controllers}
\author{Kanishk Bairagi\\Independent Researcher, India\\\texttt{kanishk.bairagi00@gmail.com}}

\begin{document}
\maketitle
\linenumbers

\begin{abstract}
We compare two LLMs, \textit{gemini-3.6-flash} and \textit{gpt-oss-120b} (via Groq), on generating Jest unit tests for Express.js controllers, using repeated trials rather than a single run. Set~A contains 25 controllers written by the author (5 trials per model per controller, 250 suites). Set~B contains 25 controllers from 13 public repositories, chosen by a rule-based protocol (3 trials, 150 suites). A deterministic non-LLM generator serves as a descriptive baseline. On Set~A, 111 of 125 Gemini suites and 99 of 125 Groq suites passed every test (88.8\% vs.\ 79.2\%; the confidence intervals overlap). On the 22 controllers that executed in all trials for both models, Gemini achieved higher coverage (line $+6.6$ points, 95\% CI [5.5, 7.8]) and a higher mutation score ($+13.5$ points, [10.7, 16.8]), but 9 of 25 (Gemini) and 15 of 25 (Groq) controllers flipped between passing and failing across trials. On Set~B, only 8 of 75 Gemini and 5 of 75 Groq suites passed every test, and the registered comparison found no detectable difference (line coverage $-6.2$ points, [$-14.8$, 4.6]). Many Set~B failures were crashes before any test ran (missing packages, wrong relative import paths). A post hoc analysis restricted to trials where both models' tests ran favored Gemini (line $+17.5$ points, [3.2, 31.4]). Results are specific to these two models, this prompt, and this harness.
\end{abstract}

\smallskip
\noindent\textbf{Keywords:} Large language models, Automated unit test generation, Express.js, Mutation testing, Empirical software engineering, Repeated trials

% ============================================================
\section{Introduction}
Writing unit tests for HTTP handlers is repetitive, and large language models (LLMs) are a natural candidate for automating it~\cite{wang2024survey}. Most evaluations of LLM test generation target Java or Python, and many report one sample per task~\cite{yuan2023chatgpt,siddiq2024junit,yang2024evaluation}. This matters because LLM outputs vary from run to run~\cite{ouyang2023nondeterminism}, and because test generators are often evaluated on code chosen or written by the evaluator.

We study Express.js controllers written as native ES Modules and tested with Jest. We ask:
\begin{itemize}
\item \textbf{RQ1.} On author-written controllers (Set~A), how do the two models compare in suite executability, coverage, and mutation score?
\item \textbf{RQ2.} How much do outcomes vary across repeated trials of the same prompt?
\item \textbf{RQ3.} Do the Set~A findings hold on external code (Set~B)?
\item \textbf{RQ4.} How do model-generated suites compare with a naive baseline, and what are the main failure causes?
\end{itemize}

\textbf{Contributions.} (i) A repeated-trial harness and the complete trial records (prompts, raw responses, Jest and Stryker outputs); (ii) a two-dataset comparison in which the second dataset is selected by a protocol committed before generation; (iii) a failure analysis showing that environment and path problems, not only test logic, dominate on external code.

\textbf{Relation to an earlier single-run draft.} An earlier draft reported a single run per controller. Re-examination of that pilot found errors (a coverage-parsing bug that scored three suites as crashes, a prompt description that did not match the code, and mutation figures for one model without a corresponding run). That draft is superseded; the pilot is reported in a corrected form in Section~\ref{sec:pilot}.

% ============================================================
\section{Related Work}

\textbf{LLM-based unit test generation and evaluation.}
Recent research has increasingly explored large language models for automated unit test generation~\cite{wang2024survey}. Early neural approaches framed test generation as translation from focal methods using sequence-to-sequence transformers~\cite{tufano2020transformers}. Modern foundation models generate entire test suites directly from source code and docstrings. Yuan et al.~\cite{yuan2023chatgpt} evaluated ChatGPT on Java unit test generation across 100 projects, observing that while generated tests exhibit high readability, many raw outputs fail compilation or runtime execution due to hallucinated identifiers and incorrect assertions. To mitigate execution failures, ChatUniTest~\cite{xie2023chatunitest} introduced an iterative generation-validation-repair loop. In an empirical evaluation of LLMs on unit test generation, Yang et al.~\cite{yang2024evaluation} showed that high line coverage does not directly translate to high fault detection, underscoring the necessity of multi-dimensional evaluation encompassing suite executability, Istanbul-style coverage, and mutation testing. Mutation-guided prompting has also been proposed to strengthen test oracle quality~\cite{dakhel2024mutap}, and CodaMOSA combines search-based testing with LLM suggestions to escape coverage plateaus~\cite{lemieux2023codamosa}. In industrial practice, Meta deployed automated unit test improvement using LLMs with continuous test execution feedback~\cite{alshahwan2024meta}.

\textbf{JavaScript and web application test generation.}
Most empirical benchmarks evaluate Java or Python programs~\cite{fraser2011evosuite,lukasczyk2022pynguin,siddiq2024junit}. In the JavaScript ecosystem, classical tools such as JSEFT focused on client-side browser DOM manipulation~\cite{mirshokraie2015jseft}. More recently, TestPilot~\cite{schafer2024testpilot} evaluated LLMs on generating tests for JavaScript npm packages, demonstrating that prompting models with runtime feedback significantly improves test quality. However, TestPilot targeted standalone library APIs with deterministic input-output functions. Testing backend HTTP route handlers and web application controllers introduces distinct architectural challenges. As Arcuri~\cite{arcuri2019evomaster} observed in the broader context of RESTful API testing, web endpoints operate through asynchronous request/response pipelines, HTTP status codes, query/body serialization, and implicit middleware chains. Evaluating unit test generation on Express.js controllers requires validating whether models can correctly navigate framework-specific conventions.

\textbf{Mocking and dependency isolation.}
Unit testing of web controllers inherently requires isolating the handler under test from external databases, authentication tokens, and third-party services. In an empirical study of mocking practices across open-source systems, Spadini et al.~\cite{spadini2017mock} demonstrated that developers predominantly use mock objects at architectural boundary layers (controllers and external interfaces) to simulate complex dependencies. In automated test generation, Fraser and Arcuri~\cite{fraser2011evosuite} noted that lack of environmental mocking prevents test execution or causes unwanted side effects. When LLMs generate unit tests without visibility into external services, mocking decisions become a primary source of fragility: models must either invent mock implementations or assume the presence of auxiliary testing utilities.

\textbf{The execution barrier and environment misalignment.}
A recurring finding in LLM code generation is that generated tests often encounter an ``execution barrier'' prior to running any test assertion. Empirical studies by Yang et al.~\cite{yang2024evaluation} and Schäfer et al.~\cite{schafer2024testpilot} reveal that execution failures in generated suites are frequently caused by unresolvable imports, missing package dependencies, and invalid module paths rather than subtle semantic defects. When language models are prompted with isolated source files without access to package manifests or filesystem layouts, models frequently make assumptions about module locations or uninstalled testing libraries. Analyzing this failure mode separately from test logic failure is essential for realistic benchmark interpretation.

\textbf{Stochastic replication and baseline calibration.}
Non-determinism is an intrinsic property of autoregressive LLM sampling~\cite{ouyang2023nondeterminism}. Ouyang et al.~\cite{ouyang2023nondeterminism} showed that identical prompts produce substantial output variability across trials. For stochastic test generation algorithms, statistical guidelines by Arcuri and Briand~\cite{arcuri2014hitchhiker} mandate repeated independent trials evaluated via non-parametric paired tests and confidence intervals. Furthermore, foundational work on automated testing baselines~\cite{arcuri2010random} emphasizes that simple random or structural invocation can achieve non-trivial code coverage simply by entering method bodies. Evaluating LLM generators without a baseline risks attributing basic structural execution to model reasoning. Coverage alone correlates weakly with fault detection~\cite{inozemtseva2014coverage}, which motivates mutation testing~\cite{jia2011mutation}.

\textbf{Positioning of this study.}
Our work connects these methodological themes. Specifically, this study provides:
(i) a repeated-trial protocol ($k=5$ and $k=3$) with exact Wilcoxon signed-rank and cluster-bootstrap inference on Express.js native ES Modules;
(ii) an evaluation across two complementary datasets, including an external 25-controller suite selected via a prespecified protocol from public repositories;
(iii) an integrated multi-metric evaluation combining execution success, Istanbul line, branch, and function coverage, and StrykerJS mutation testing;
(iv) a deterministic non-LLM baseline to calibrate the empirical floor of naive handler execution;
(v) an empirical taxonomy of environment-induced import crashes versus assertion failures; and
(vi) replication artifacts for all prompts, responses, execution logs, and analysis scripts.

% ============================================================
\section{Study Design}
\subsection{Datasets}
\textbf{Set~A} consists of 25 controllers written by the author, covering domains such as authentication, payments, uploads, and health checks (all native ES Modules). Because the author also designed the prompt and the harness, Set~A may favor either model; Set~B exists to test this.

\textbf{Set~B} consists of 25 controller or route-handler files from 13 public GitHub repositories (Table~\ref{tab:external_manifest}). The selection protocol (\texttt{external\_selection\_protocol.md}) was committed on 2026-10-04 at 21:51 (UTC+5:30), before the candidate search ran; the manifest was committed at 22:50. The protocol required: GitHub search for Express.js backends with at least 50 stars; a permissive license (MIT, Apache-2.0, BSD), checked from the LICENSE file; a pinned commit; not authored by the study author; files of 40--300 lines with at least three exported handlers; a local import closure of at most five files; no import-time side effects. Within each repository, qualifying files were taken in alphabetical path order, at most three per repository, until 25 were selected; every candidate and exclusion reason was recorded. One selected file crashed at import time (a Stripe client created without a key); it was excluded under the protocol, and the next qualifying file by the same rule replaced it. Files were copied unchanged with their import closures, and their dependencies resolve against the study project's own \texttt{package.json}, not the originals. All 25 are written with ES module syntax. Several repositories are popular and may appear in the models' training data (Table~\ref{tab:external_manifest}).

\subsection{Models, prompt, and settings}
We used \textit{gemini-3.6-flash} through the Google Gemini API (\texttt{@google/genai} 2.22.0, \texttt{generateContent})~\cite{google2026gemini36} and \textit{gpt-oss-120b} through Groq's OpenAI-compatible endpoint (\texttt{openai} client 7.15.0)~\cite{openai2025gptoss}. Both models received the same system prompt and user message (below). Sampling parameters, thinking effort, and reasoning effort were left at provider defaults; the output-token cap was set to 32{,}768 for both. Google's documentation recommends leaving Gemini~3 sampling parameters at their defaults~\cite{google2026gemini36}. Gemini reports only its model alias in responses; Groq reported 41 distinct backend fingerprints across the 125 Set~A calls, so Groq's backend was not fixed during the experiment. Free-tier access was used for both providers.

\begin{lstlisting}
System prompt:
You are a Senior QA Automation Engineer. Output only valid Jest unit tests using ES Module syntax. Always import target functions using relative path '../dataset/${targetFilename}'. Return ONLY valid JavaScript code in a clean markdown block (```javascript ... ```). Do NOT include any conversational filler, explanations, markdown text outside the code block, or questions.

User message:
Target controller: ../dataset/${targetFilename}

Code to test:
${sourceCode}
\end{lstlisting}

The prompt does not tell the model where the generated test file will be stored, which packages are installed, or that Jest runs in ESM mode.

\subsection{Harness}
For each (model, controller, trial) the harness (Node v24.21.0; Jest 30.5.1~\cite{jest2026}; Stryker 10.0.0~\cite{stryker2026}) requests one completion, extracts the code block, and applies a sanitizer that (a) injects \texttt{import \{ jest \} from '@jest/globals'} when the code uses \texttt{jest.} and the string \texttt{@jest/globals} is absent, and (b) rewrites imports of the target controller to \texttt{../dataset/<file>.js}. The sanitizer does not detect the case where \texttt{@jest/globals} is imported without \texttt{jest}. The suite is written to a flat test directory and run once with Jest in ESM mode (\texttt{--experimental-vm-modules}, \texttt{--coverage}, 35-second timeout), parsing the JSON report. Trials run in trial-major order (trial 1 for all controllers, then trial 2, \dots). A request is repeated only when the provider returned no output: after transport-level errors or HTTP 503 responses (up to five attempts, waiting 3, 6, 9, and 12 seconds), after per-minute rate-limit responses (10-second waits, up to three times), and, if all attempts fail, in a later session. Outputs are never re-sampled because of their content: truncated, empty, or unparseable outputs are recorded as results. The 32{,}768-token cap was not reached by any call (all finish reasons were normal stops).

\textbf{Operational details.} The Gemini free tier allowed 20 requests per day, so Set~A was spread over several days and used 9 Gemini API keys from 9 accounts (billing not enabled); the harness does not record which key served which call. Set~B used 20 Gemini API keys from 20 accounts (generation dates: 2026-10-05 to 2026-10-07). TLS certificate verification was disabled for API traffic because of a local network proxy.

\subsection{Metrics}
\textit{Executable}: Jest produced a report for the suite. \textit{Fully passing}: all tests passed. Coverage (statement, line, branch, function) is taken from Jest's Istanbul coverage map for the target file~\cite{istanbul2026}; line coverage is computed from the statement map. \textit{Mutation score} is computed by StrykerJS~\cite{stryker2026} only for suites in which every test passed (Stryker requires a passing initial run); it is reported as \emph{strict} (killed/total) and \emph{Stryker} ((killed+timeout)/total). Suites that cannot be mutated receive no score, and no zeros are imputed.

\subsection{Statistical analysis and deviations}
The unit of analysis is the controller: trial values are averaged within a controller, and models are compared with a two-sided exact Wilcoxon signed-rank test on per-controller paired differences, with 95\% percentile bootstrap confidence intervals (10{,}000 resamples, seed 42) that resample controllers (Set~A) or repositories (Set~B, because up to three files share a repository). Holm correction is applied across the five primary outcomes (line, branch, function coverage, two mutation scores)~\cite{holm1979,wilcoxon1945,efron1979}. With $n=22$ and all differences in one direction, the exact p-value reaches its minimum, so we emphasize effect sizes and intervals.

An analysis plan was drafted before any trial was analyzed, but it was first committed to the repository on 2026-10-04 at 23:20 (UTC+5:30), after the Set~A analysis had been run and committed; we therefore do not claim pre-registration for Set~A. The Set~B selection protocol was committed on 2026-10-04 at 21:51, before the Set~B candidate search, and the Set~B addendum to the plan was committed at 23:20, before Set~B generation began (first request at about 23:35). Deviations from the plan:
\begin{enumerate}
\item The plan specified trial-level matching. The primary Set~A comparison instead used the 22 controllers that executed in all five trials for both models; trial-level matching (sensitivity A) and an all-controllers analysis with zeros for non-executable suites (sensitivity B, secondary) are reported alongside.
\item In Set~B, ``executable'' counted suites that crashed before running any test. Sensitivity C (post hoc) restricts to trials where both models ran at least one test.
\item Set~B mutation testing remained limited by the rule that only fully passing suites are mutated.
\end{enumerate}

% ============================================================
\section{Results}
\subsection{Pilot run}\label{sec:pilot}
The pilot ran each controller once per model, across several sessions on two days, using the same prompt. After correcting a coverage-parsing bug, Gemini produced 25 of 25 executable and fully passing suites (439 tests), while Groq produced 21 of 25 executable and 17 of 25 fully passing suites (244 passed, 4 failed tests); four Groq outputs were cut off mid-statement. Output-token limits were not set in the pilot. Table~\ref{tab:pilot_run_corrected} gives the corrected values.
\tablein{table1_pilot_run.tex}

\subsection{Set A: author-written controllers (RQ1, RQ2)}
% SRC: results/analysis/suite_success_summary.csv, per_trial_summary.csv, primary_controller_comparison.csv
Across five trials, 111 of 125 Gemini suites (88.8\%, 95\% CI [81.6, 95.2]) and 99 of 125 Groq suites (79.2\%, [69.6, 87.2]) passed every test. Gemini produced 12 assertion-failure suites and 2 syntax errors; Groq produced 25 assertion-failure suites and 1 syntax error. No output was truncated.

Outcomes varied between trials: 9 of 25 controllers flipped between fully passing and failing for Gemini, and 15 of 25 for Groq (Table~\ref{tab:seta_suite_success}). Per-trial means across the five trials were 24.6 (SD 0.55) executable and 22.2 (SD 1.64) fully passing suites for Gemini, and 24.8 (0.45) and 19.8 (1.30) for Groq.
\tablein{table3_seta_suite_success.tex}

On the 22 controllers that executed in all trials for both models (excluded: \texttt{02\_product}, \texttt{10\_analytics}, \texttt{22\_twofactor}), Gemini's mean line coverage was 94.1\% vs.\ 87.5\% for Groq (difference $+6.6$ points, 95\% CI [5.5, 7.8]); branch coverage was higher by 7.9 points [5.5, 10.6], function coverage by 18.4 points [12.7, 24.3], and the strict mutation score by 13.5 points [10.7, 16.8] (Table~\ref{tab:seta_comparison}). Trial-level matching (25 controllers) gave nearly identical differences, and the all-controller analysis with zeros gave smaller but still positive differences. Function coverage was low for both models (61.4\% and 43.0\%) even when line coverage was high. Mutation scores describe only fully passing suites (111 Gemini and 99 Groq controller-trials), so they measure the strength of passing tests, not overall suite quality.
\tablein{table2_seta_comparison.tex}

\textbf{Cost.} Mean output tokens per call were about 4{,}726 for Gemini (3{,}044 completion plus 1{,}682 thinking tokens) and 2{,}881 for Groq (2{,}463 completion plus 418 reasoning tokens; Table~\ref{tab:token_usage}); Gemini's higher coverage should be read together with this larger generation budget and with its larger tests (426 vs.\ 295 tests passed per trial).

\subsection{Set B: external controllers (RQ3)}
% SRC: results/analysis_external/*.csv
Only 8 of 75 Gemini suites (10.7\%, repository-cluster 95\% CI [1.4, 24.2]) and 5 of 75 Groq suites (6.7\%, [0.0, 17.4]) passed every test. In the registered primary comparison (21 controllers, three trials each; Table~\ref{tab:setb_comparison}), Gemini's line coverage was lower by 6.2 points (95\% CI [$-14.8$, 4.6]); branch ($-1.9$ points) and function ($-4.8$ points) differences also had intervals containing zero, and the trial-level and zero-imputed analyses agreed. Only one controller had fully passing suites for both models, so a mutation comparison was not possible.

Of the 150 Set~B trials, 62 ran no test at all (42 Gemini, 20 Groq), almost all because the test file crashed on import. A post hoc classification of the 137 non-passing trials (Table~\ref{tab:setb_failures}) shows: 27 trials imported a package that is not installed (\texttt{supertest}; 17 Gemini, 10 Groq); 28 imported a relative path that is valid from the controller's directory but not from the test file's directory (24 Gemini, 4 Groq); 1 Groq trial invented a path; 5 were syntax errors; and 75 were assertion failures (25 Gemini, 50 Groq). The first two categories come from the prompt not specifying the test location or installed packages, so they reflect the harness as much as the models.

\textbf{Post hoc sensitivity C} (exploratory) restricts to controller-trials where both models' tests ran at least one test (19 controllers from 11 repositories, 31 trial pairs). Here Gemini's coverage was higher: line $+17.5$ points [3.2, 31.4], branch $+24.4$ [7.3, 41.5], function $+20.2$ [4.2, 36.3] (Holm-adjusted $p=0.040$ for each). This analysis conditions on an outcome (tests running) and uses a subset, so it shows a pattern, not an unbiased estimate.
\tablein{table4_setb_comparison.tex}
\tablein{table5_setb_failures.tex}

\subsection{Naive baseline (RQ4)}
% SRC: results/baseline_naive/*.csv, results/analysis/*.csv, results/analysis_external/*.csv
The baseline is a deterministic script that, for each exported handler, makes three calls with mocked request, response, and next objects and asserts only that the call settles. It uses no LLM and no controller-specific edits. On the 22 matched Set~A controllers its mean line, branch, and function coverage was 37.8\%, 33.9\%, and 39.5\%, against 94.1\%, 94.6\%, and 61.4\% (Gemini) and 87.5\%, 86.8\%, and 43.0\% (Groq). On Set~B (19-controller subset) the baseline's mean line coverage was 34.3\% over all 19 controllers (43.4\% over the 15 whose baseline suite ran), against 50.8\% (Gemini) and 33.3\% (Groq); the baseline's function coverage (57.0\%, or 72.2\% over runnable suites) exceeded both models' (50.3\% and 30.2\%) because it calls every exported function (Table~\ref{tab:naive_baseline}). These comparisons are descriptive: the baseline is a single deterministic run, the model figures are means over matched trials, and the populations of suites differ. The baseline's only assertion is that a call settles, so it cannot detect faults; we did not mutation-test it.
\tablein{table6_naive_baseline.tex}
\tablein{table7_token_usage.tex}

% ============================================================
\section{Discussion}
\textbf{Single runs are unreliable here.} The pilot suggested that Gemini was far ahead (25 of 25 suites passing vs.\ 17 of 25 for Groq). With five trials the gap in fully passing suites is smaller (88.8\% vs.\ 79.2\%, overlapping intervals), and 36\% of controllers for Gemini and 60\% for Groq changed outcome between trials. A single sample of each model would have misrepresented both the size of the difference and the stability of the results.

\textbf{Author-written versus external code.} On Set~A both models far exceeded the baseline on line and branch coverage, and Gemini led Groq on all four measures. On Set~B the registered comparison found no detectable difference, few suites passed, and the baseline matched or exceeded the models on several measures. One plausible explanation, which we did not test, is that models write tests that depend on project structure and packages they cannot see. Both the dataset and the harness may contribute: Set~A was written by the same author who designed the prompt, and Set~B's prompt gave no information about test location or installed dependencies. A prompt variant that states these facts is the obvious next experiment.

\textbf{Gemini versus Groq.} When tests ran, Gemini's tests usually covered more code and were stronger under mutation (Set~A), but Gemini's Set~B suites more often crashed before running (42 vs.\ 20 of 75). We cannot say whether this reflects the model's habits (for example, relative import style) or an interaction with the harness.

% ============================================================
\section{Threats to Validity}
\textbf{Construct validity.} Coverage and mutation score are proxies for test quality. Function coverage was low even when line coverage was high. Mutation scores exist only for fully passing suites, which excludes the suites most likely to be weak, and Set~B mutation data are too sparse to compare.

\textbf{Internal validity.} The harness is part of the treatment: the sanitizer, the flat test directory, the Jest timeout, and the prompt wording all affect which suites run. The prompt does not state the test location or installed packages, which caused many Set~B crashes. Free-tier access, HTTP 503 responses, multiple Gemini keys, and Groq's changing backend fingerprints mean the services were not controlled; the model alias for Gemini does not guarantee an identical model across days. Thinking and reasoning effort were left at provider defaults, and the two providers count reasoning tokens differently. TLS verification was disabled. Set~A, Set~B, and the baseline were run in different sessions.

\textbf{External validity.} Two models, one framework, one language, one prompt, and no prompt tuning. Set~B has 25 files from 13 repositories, selected by a protocol but still a small, search-dependent sample. Public repositories may be in the models' training data. Set~A is author-written.

\textbf{Conclusion validity.} Clusters are few (22 controllers in Set~A, 11 to 13 repositories in Set~B), so intervals are wide. Several analyses are post hoc and labeled as such; we report deviations from the plan above. We did not compare against a paid frontier model or a human-written test suite.

% ============================================================
\section{Conclusion}
In a repeated-trial comparison, Gemini's generated Jest suites outperformed Groq's on author-written Express.js controllers, with moderate but consistent coverage and mutation-score advantages, and both models changed outcomes between identical trials. On external controllers few suites passed, the registered comparison showed no detectable difference, and a trivial script reached coverage comparable to the models on several measures. Single-run results and author-written benchmarks can overstate the usefulness of LLM test generation; prompts that specify test location and dependencies, and repair loops, are the next experiments.

\section*{Data and code availability}
Harness, datasets, selection protocol and manifest, all raw trial records (prompts, responses, Jest and Stryker outputs), analysis scripts, and the lab notebook are available in the public repository at \url{https://github.com/kanishkbairagi/express_llm_benchmark1}, organized by milestone tags corresponding to the study progression (\texttt{trials-end}, \texttt{mutation-end}, \texttt{analysis-v2}, \texttt{setb-end}, \texttt{setb-analysis-v2}, \texttt{baseline-v1}, \texttt{paper-tables-v1}). API keys are not included. External files retain their original licenses.

\section*{Declarations}
\textbf{Funding:} none. \textbf{Conflicts of interest:} none declared. \textbf{Ethics:} not applicable.

% ============================================================
\begin{thebibliography}{99}
\small
\bibitem{wang2024survey} J.~Wang, Y.~Huang, C.~Chen, Z.~Liu, S.~Wang, Q.~Wang. Software testing with large language models: Survey, landscape, and vision. \textit{IEEE Trans.\ Softw.\ Eng.}, 50(4):803--836, 2024.
\bibitem{yuan2023chatgpt} Z.~Yuan, M.~Liu, S.~Ding, K.~Wang, Y.~Chen, X.~Peng, Y.~Lou. Evaluating and improving ChatGPT for unit test generation. In \textit{Proc.\ ACM Softw.\ Eng.} (FSE), 1:1703--1726, 2024.
\bibitem{siddiq2024junit} M.~L.~Siddiq et al. Using large language models to generate JUnit tests: An empirical study. In \textit{Proc.\ EASE}, pp.\ 242--252, 2024.
\bibitem{yang2024evaluation} L.~Yang et al. On the evaluation of large language models in unit test generation. In \textit{Proc.\ ASE}, pp.\ 581--593, 2024.
\bibitem{ouyang2023nondeterminism} S.~Ouyang, J.~M.~Zhang, M.~Harman, M.~Wang. LLM is like a box of chocolates: The non-determinism of ChatGPT in code generation. arXiv:2308.02828, 2023.
\bibitem{lemieux2023codamosa} C.~Lemieux, J.~P.~Inala, S.~K.~Lahiri, S.~Sen. CodaMOSA: Escaping coverage plateaus in test generation with pre-trained large language models. In \textit{Proc.\ ICSE}, pp.\ 919--931, 2023.
\bibitem{schafer2024testpilot} M.~Sch\"afer, S.~Nadi, A.~Eghbali, F.~Tip. An empirical evaluation of using large language models for automated unit test generation. \textit{IEEE Trans.\ Softw.\ Eng.}, 50(1):85--105, 2024.
\bibitem{xie2023chatunitest} Z.~Xie, Y.~Chen, C.~Zhi, S.~Deng, J.~Yin. ChatUniTest: A ChatGPT-based automated unit test generation tool. arXiv:2305.04764, 2023.
\bibitem{alshahwan2024meta} N.~Alshahwan et al. Automated unit test improvement using large language models at Meta. In \textit{Companion Proc.\ FSE}, pp.\ 497--501, 2024.
\bibitem{dakhel2024mutap} A.~M.~Dakhel et al. Effective test generation using pre-trained large language models and mutation testing. \textit{Inf.\ Softw.\ Technol.}, 171:107456, 2024.
\bibitem{tufano2020transformers} M.~Tufano, D.~Drain, A.~Svyatkovskiy, S.~K.~Deng, N.~Sundaresan. Unit test case generation with transformers and focal context. arXiv:2009.05617, 2020.
\bibitem{fraser2011evosuite} G.~Fraser, A.~Arcuri. EvoSuite: Automatic test suite generation for object-oriented software. In \textit{Proc.\ ESEC/FSE}, pp.\ 415--418, 2011.
\bibitem{lukasczyk2022pynguin} S.~Lukasczyk, G.~Fraser. Pynguin: Automated unit test generation for Python. In \textit{Companion Proc.\ ICSE}, pp.\ 165--169, 2022.
\bibitem{mirshokraie2015jseft} S.~Mirshokraie, A.~Mesbah, K.~Pattabiraman. JSEFT: Automated JavaScript unit test generation. In \textit{Proc.\ ICST}, pp.\ 1--10, 2015.
\bibitem{arcuri2019evomaster} A.~Arcuri. RESTful API automated test case generation with EvoMaster. \textit{ACM Trans.\ Softw.\ Eng.\ Methodol.}, 28(1):3:1--3:37, 2019.
\bibitem{spadini2017mock} D.~Spadini, M.~Aniche, M.~Bruntink, A.~Bacchelli. To mock or not to mock? An empirical study on mocking practices. In \textit{Proc.\ MSR}, pp.\ 402--412, 2017.
\bibitem{arcuri2010random} A.~Arcuri, M.~Z.~Iqbal, L.~Briand. Formal analysis of the effectiveness and predictability of random testing. In \textit{Proc.\ ISSTA}, pp.\ 219--230, 2010.
\bibitem{inozemtseva2014coverage} L.~Inozemtseva, R.~Holmes. Coverage is not strongly correlated with test suite effectiveness. In \textit{Proc.\ ICSE}, pp.\ 435--445, 2014.
\bibitem{jia2011mutation} Y.~Jia, M.~Harman. An analysis and survey of the development of mutation testing. \textit{IEEE Trans.\ Softw.\ Eng.}, 37(5):649--678, 2011.
\bibitem{arcuri2014hitchhiker} A.~Arcuri, L.~Briand. A hitchhiker's guide to statistical tests for assessing randomized algorithms in software engineering. \textit{Softw.\ Test.\ Verif.\ Reliab.}, 24(3):219--250, 2014.
\bibitem{holm1979} S.~Holm. A simple sequentially rejective multiple test procedure. \textit{Scand.\ J.\ Statist.}, 6(2):65--70, 1979.
\bibitem{wilcoxon1945} F.~Wilcoxon. Individual comparisons by ranking methods. \textit{Biometrics Bulletin}, 1(6):80--83, 1945.
\bibitem{efron1979} B.~Efron. Bootstrap methods: Another look at the jackknife. \textit{Ann.\ Statist.}, 7(1):1--26, 1979.
\bibitem{openai2025gptoss} OpenAI. Introducing gpt-oss. OpenAI Research, August 2025. \url{https://openai.com/index/introducing-gpt-oss/}.
\bibitem{google2026gemini36} Google. Gemini models: Gemini 3.6 Flash. Google AI for Developers, 2026. \url{https://ai.google.dev/gemini-api/docs/models/gemini}.
\bibitem{jest2026} OpenJS Foundation. Jest: Delightful JavaScript testing. Version 30.5.1, 2026. \url{https://jestjs.io/}.
\bibitem{istanbul2026} Istanbul.js contributors. Istanbul: JavaScript code coverage tool. 2026. \url{https://istanbul.js.org/}.
\bibitem{stryker2026} Stryker Mutator Team. StrykerJS: Mutation testing for JavaScript and TypeScript. Version 10.0.0, 2026. \url{https://stryker-mutator.io/}.
\end{thebibliography}

\appendix
\section{External files used in Set~B}
\tablein{table8_external_manifest.tex}
\end{document}
